
% ==========================================
% 1. FONTS, LINGUA E FORMATTAZIONE BASE
% ==========================================
\usepackage[T1]{fontenc}
\renewcommand*\familydefault{\sfdefault} % Font sans-serif
\usepackage{setspace}
\usepackage{multicol}
\usepackage{relsize}
\usepackage[normalem]{ulem}

% ==========================================
% 2. MATEMATICA E SIMBOLI
% ==========================================
\usepackage{amsmath, amssymb, mathtools}
\usepackage{upgreek}
\usepackage{nicefrac}
\usepackage{cancel}
\usepackage{wasysym}
\usepackage{pifont}
\usepackage{bm}

% ==========================================
% UNITÀ DI MISURA
% ==========================================
\usepackage{siunitx}

% ==========================================
% 4. COLORI E GRAFICA
% ==========================================
\usepackage[dvipsnames,table]{xcolor}
\usepackage{graphicx}
\usepackage{svg}
\usepackage{tikz}
\usetikzlibrary{quotes,angles}
\usepackage{circuitikz}
\usepackage{pgfplots}
\pgfplotsset{compat=1.18}
\usepackage{wrapfig}
\usepackage{tcolorbox}
\usepackage{rotating}
\usepackage{chemfig}
\usepackage{soul} % Gestisce l'evidenziazione su più righe

% ==========================================
% 5. TABELLE ED ELENCHI
% ==========================================
\usepackage{tabularx, tabularray, longtable, array, booktabs, multirow}
\usepackage{enumitem}

% ==========================================
% 6. GEOMETRIA PAGINA E SPAZIATURE
% ==========================================
\usepackage{geometry}
\geometry{
    top=1cm,
    bottom=1.5cm,
    left=1cm,
    right=1cm
}
\usepackage{afterpage}
\usepackage{float}
\usepackage{needspace}
\usepackage{etoolbox}
\flushbottom

% Controllo vedove e orfani
\widowpenalty=100    
\clubpenalty=100     

% ==========================================
% 7. DATE E DIDASCALIE
% ==========================================
\usepackage[style=ddmmyyyy]{datetime2}
\DTMsetup{datesep=/}
\usepackage{caption}
\captionsetup[figure]{labelformat=empty} % Rimuove "Figure n:"

% ==========================================
% 8. NOTE A PIÈ DI PAGINA
% ==========================================
% \usepackage[para,splitrule,bottom]{footmisc}
% \renewcommand{\footnotesize}{\scriptsize} 
% \setlength{\skip\footins}{3pt plus 2pt minus 1pt} 
% \setlength{\footnotesep}{0.1em}   

\usepackage[para,splitrule,bottom]{footmisc}
\renewcommand{\footnotesize}{\scriptsize} 
\setlength{\skip\footins}{15pt plus 4pt minus 2pt} % Distanza SOPRA la linea
\setlength{\footnotesep}{0.8em}                    % Distanza SOTTO la linea


% ==========================================
% 9. GESTIONE TITOLI E GERARCHIA VISIVA
% ==========================================
\usepackage{titlesec}
\usepackage{titletoc}
\usepackage{tocloft}
\usepackage{titleref}

% Definizioni Colori Standard e Titoli
\definecolor{myred}{rgb}{0.7,0.1,0.1} % Rosso principale (\section)
\definecolor{cobalto}{RGB}{0, 71, 171} % Blu elettrico saturo (\subsection)
\definecolor{verdone}{RGB}{0, 120, 50} % Verde scuro e saturo, zero rosso (\subsubsection)
\definecolor{arancioterroso}{RGB}{214, 76, 29} % Verde scuro e 


\definecolor{CadetBlue}{RGB}{95,158,160}
\definecolor{azzurrocurve}{RGB}{116,159,211}
\definecolor{verdecurve}{RGB}{79,126,58}
\definecolor{rossocurve}{RGB}{214,32,45}
\definecolor{giallocurve}{RGB}{224,154,74}
\definecolor{gialloevidenziatore}{RGB}{223,255,0}
\definecolor{grigino}{gray}{0.35} 

% \section: Enorme, rosso, con linea separatrice sotto
\titleformat{\section}[hang]
  {\Huge\bfseries\color{myred}}
  {\thesection}{1em}{}
  [\vspace{-0.5ex}\rule{\textwidth}{0.8pt}]
\titlespacing*{\section}{0pt}{4ex plus 1ex minus .2ex}{2.5ex plus .2ex}

% \subsection: Molto grande, cobalto scuro e saturo
\titleformat{\subsection}[hang]
  {\LARGE\bfseries\color{cobalto}}
  {\thesubsection}{0.8em}{}
\titlespacing*{\subsection}{0pt}{3.5ex plus 1ex minus .2ex}{1.5ex plus .2ex}

% \subsubsection: Grande, verdone scuro (stacco netto da rosso e blu)
\titleformat{\subsubsection}[hang]
  {\Large\bfseries\color{verdone}}
  {\thesubsubsection}{0.8em}{}
\titlespacing*{\subsubsection}{0pt}{2.5ex plus 1ex minus .2ex}{1ex plus .2ex}

% \paragraph: Evidenziato, grigio scuro, a capo
\titleformat{\paragraph}[hang]
  {\large\bfseries\color{arancioterroso}}
  {\theparagraph}{1em}{}
\titlespacing*{\paragraph}{0pt}{2ex plus 0.5ex minus .2ex}{0.5ex plus .2ex}


% ── profondità indice e numerazione ──────
\setcounter{tocdepth}{5}     % mostra \paragraph nell'indice
\setcounter{secnumdepth}{5}  % numera \paragraph nel testo

\usepackage{ocgx2}

% ── GESTIONE VISIVA INDICE (tocloft) ──────
\cftpagenumberson{section}
\cftpagenumberson{subsection}
\cftpagenumberson{subsubsection}
\cftpagenumbersoff{paragraph}

% 2. Gerarchia visiva: Section (Enorme)
\renewcommand{\cftsecfont}{\Large\bfseries}
\renewcommand{\cftsecpagefont}{\Large\bfseries}
\setlength{\cftbeforesecskip}{2ex} % Spazio extra prima delle sezioni

% 3. Gerarchia visiva: Subsection (Eredita la vecchia section)
\renewcommand{\cftsubsecfont}{\Large\bfseries}
\renewcommand{\cftsubsecpagefont}{\Large\bfseries}
 \setlength{\cftbeforesubsecskip}{1.5ex} % Spazio extra prima delle sottosezioni

% 4. Gerarchia visiva: Subsubsection (Leggermente maggiorata)
\renewcommand{\cftsubsubsecfont}{\large}
\setlength{\cftbeforesubsecskip}{1ex} % Spazio extra prima delle sottosottosezioni

% 5. Compensazione larghezza numeri (Evita che "1.1.2.1" si sovrapponga al testo)
\setlength{\cftsecnumwidth}{2.5em}
\setlength{\cftsubsecnumwidth}{3.2em}
\setlength{\cftsubsubsecnumwidth}{4.2em}


% ==========================================
% 10. MACRO E AMBIENTI PERSONALIZZATI
% ==========================================
\newcommand{\pro}{{\textcolor{ForestGreen}{\textbf{V}}}: }
\newcommand{\spuntaverde}{{\textcolor{ForestGreen}{\textbf{V}}}}
\newcommand{\contro}{{\textcolor{red}{\textbf{X}}}: }
\newcommand{\crocerossa}{{\textcolor{red}{\textbf{X}}}}
\newcommand{\fd}{\ensuremath{ \rightarrow \hspace{0.5em} }}   % freccia a destra
\newcommand{\fs}{\ensuremath{ \leftarrow }}   % freccia sinistra
\newcommand{\fds}{\ensuremath{ \leftrightarrow }} % freccia destra e sinistra
\newcommand{\fdn}{\ensuremath{ \not\rightarrow }}   % freccia destra negata
\newcommand{\fsn}{\ensuremath{ \not\leftarrow }}   % freccia sinistra negata
\newcommand{\fsu}{\ensuremath{ \uparrow \hspace{0.5em} }}   % freccia su
\newcommand{\fgiu}{\ensuremath{ \downarrow \hspace{0.5em} }}   % freccia su
\newcommand{\fsugiu}{\ensuremath{ \updownarrow \hspace{0.5em} }}   % freccia su e giù
\newcommand{\mq}{\ifmmode \text{m}^2 \else m$^2$ \fi} 
\newcommand{\mc}{\ifmmode \text{m}^3 \else m$^3$ \fi} 
\newcommand{\micro}{\ensuremath{ \mu }}
\newcommand{\s}{\ifmmode \text{s}\else s\fi} 
\newcommand{\mezzo}{\ifmmode \frac{1}{2} \else $\frac{1}{2}$\fi} 
\newcommand{\medsize}{\fontsize{12}{14}\selectfont}
\newcommand{\patm}{\ifmmode p_{\text{atm}} \else $p_{\text{atm}}$ \fi} 
\newcommand{\diam}{\ifmmode \diameter \else $\diameter$ \fi} 
\newcommand{\perpendicolare}{\ifmmode \perp \else $\perp$ \fi} 
\newcommand{\Emec}{\ifmmode E_{\text{mec}} \else $E_{\text{mec}}$ \fi} 
\newcommand{\Eth}{\ifmmode E_{\text{th}} \else $E_{\text{th}}$ \fi} 
\newcommand{\pompaprot}{\ifmmode H^+-\textnormal{-ATPasi} \else $H^+$-ATPasi \fi} 



\newcommand{\romboelenco}{\textcolor{RubineRed}{\(\;\){\textbf{\ding{117}}}\(\;\)}} 



\newcommand{\evidenzia}[1]{\colorbox{gialloevidenziatore}{\hspace{-0.1em}#1\hspace{-0.1em}}}

\newcommand{\evidenziaarancione}[1]{\colorbox{arancioevidenza}{\hspace{-0.1em}#1\hspace{-0.1em}}}

\newcommand{\itsmall}[1]{%
  \ifmmode \text{\textcolor{grigino}{\textit{\smaller({#1})}}}%
  \else \textcolor{grigino}{\textit{\smaller({#1})}} \fi}

\newcommand{\cobalto} [1] {\textcolor {cobalto}{#1}}

% Elenchi ottimizzati
\newcounter{puntidepth}
\setcounter{puntidepth}{0}

\newenvironment{punti}{%
  \stepcounter{puntidepth}%
  \ifnum\value{puntidepth}=1 \def\puntilab{{\large\textbf{-}}}%
  \else\ifnum\value{puntidepth}=2 \def\puntilab{\raisebox{0.25ex}{\scalebox{0.6}{$\blacktriangleright$}}}%
  \else\ifnum\value{puntidepth}=3 \def\puntilab{$\vcenter{\hbox{\scalebox{0.8}{$\bullet$}}}$}%
  \else\ifnum\value{puntidepth}=4 \def\puntilab{\raisebox{0.1ex}{\scalebox{0.7}{$\ast$}}}%
  \else \def\puntilab{\raisebox{0.1ex}{\scalebox{0.6}{$\triangleright$}}}%
  \fi\fi\fi\fi
  \begin{itemize}[leftmargin=1.5em, labelsep=0.3em, itemsep=-0.35em, label=\puntilab]%
}{%
  \end{itemize}%
  \addtocounter{puntidepth}{-1}%
}


\newenvironment{numeri}{\begin{enumerate}[leftmargin=1.5em, labelsep=0.3em, itemsep=-0.35em, label=\arabic*)]}{\end{enumerate}}


\newcommand{\nota}[1]{%
  \begin{tcolorbox}[colback=yellow!10, colframe=orange!80!black, title=\textbf{Nota}]
    #1
  \end{tcolorbox}%
}


\newcommand{\immagineinlinea}[3][]{%
  % #1 = didascalia opzionale (testo sotto l'immagine)
  % #2 = larghezza (es. 3cm)
  % #3 = nome del file (senza estensione)
  \raisebox{-.5\height}{%
    \parbox[t]{#2}{%
      \includegraphics[width=#2]{#3}%
      \ifx&#1&\else
        \par\vspace{2pt}\noindent\footnotesize \centering #1\par
      \fi
    }%
  }%
}

\newcommand{\immaginecentrata}[3][]{%
  % #1 = didascalia opzionale (testo sotto l'immagine)
  % #2 = larghezza (es. 3cm)
  % #3 = nome del file (senza estensione)
  \begin{center}
  \raisebox{-.5\height}{%
    \parbox[t]{#2}{%
      \includegraphics[width=#2]{#3}%
      \ifx&#1&\else
        \par\vspace{2pt}\noindent\footnotesize \centering  #1\par
      \fi
    }%
  }%
  \end{center}
}


% Definisci il colore e dillo al pacchetto soul
\definecolor{arancioevidenza}{RGB}{255, 128, 0}
\sethlcolor{arancioevidenza}

% --- COMANDI SPECIALI (Evidenziati a pennarello con testo nero) ---
% L'ordine dei \protect è chirurgico per sopravvivere all'indice

\newcommand{\sectionsp}[1]{%
  \section[\texorpdfstring{\protect\textcolor{black}{\protect\hl{#1}}}{#1}]%
          {\texorpdfstring{\protect\textcolor{black}{\protect\hl{#1}}}{#1}}%
}

\newcommand{\subsectionsp}[1]{%
  \subsection[\texorpdfstring{\protect\textcolor{black}{\protect\hl{#1}}}{#1}]%
             {\texorpdfstring{\protect\textcolor{black}{\protect\hl{#1}}}{#1}}%
}

\newcommand{\subsubsectionsp}[1]{%
  \subsubsection[\texorpdfstring{\protect\textcolor{black}{\protect\hl{#1}}}{#1}]%
                {\texorpdfstring{\protect\textcolor{black}{\protect\hl{#1}}}{#1}}%
}

\newcommand{\paragraphsp}[1]{%
  \paragraph[\texorpdfstring{\protect\textcolor{black}{\protect\hl{#1}}}{#1}]%
            {\texorpdfstring{\protect\textcolor{black}{\protect\hl{#1}}}{#1}}%
}

% ==========================================
% 11. HYPERREF (SEMPRE PER ULTIMO)
% ==========================================
\usepackage{hyperref}
\hypersetup{
    colorlinks=true,
    linkcolor=black, % Colore dei link interni
    urlcolor=blue,   % Colore dei link esterni
    citecolor=black  % Colore delle citazioni
}
